\documentclass{situationist-preprint}

\title{Galileo's Mathematization of Nature Revisited}
\msid{PHE-17-b}
\msfamily{PHE-17}
\msdiscipline{Phenomenology}
\mstype{Companion manuscript}
\mskeywords{Galileo; mathematization; life-world; Husserl; Crisis; primary and secondary qualities; scientific realism; replication}

\begin{document}
\maketitle

\begin{abstract}
Husserl called Galileo at once a discovering and a concealing genius. We separate the two achievements. The discovery, the mathematization of motion, was a founded extension of the life-world in the sense of PHE-17-a, and its evidence was given in life-world measurements, as the inclined-plane experiments show; nothing in it conflicts with realism about its results. The concealment was a separate step, the redescription of the life-world as a subjective effect of the mathematical world, made through the doctrine of secondary qualities. Because the two steps are separable, the Galilean pattern can recur in any successful science without any defect in its results, and it can be undone, as the recent correction of experimental psychology shows, without abandoning them.
\end{abstract}

\section{Two steps}

Husserl's analysis of Galileo in the \emph{Crisis} has two parts \citep[§9]{Husserl1954}. The first describes the idealization of shapes into geometrical limit-shapes and the indirect mathematization of the sensible qualities, which made it possible to calculate the course of experience. The second describes the \enquote{surreptitious substitution} of the mathematically idealized world for the world of experience, which Husserl locates in the same figure. PHE-17-a accepts the analysis and qualifies it: the substitution took centuries to complete, and the Galilean episode is the beginning of a crisis, not a crisis already present. We argue that the two parts are separable in Galileo's own work, and that the separation clarifies the recurrence thesis.

\begin{proposition}\label{p:founded}
Galileo's law of free fall is a founded extension of the life-world in the sense of Definition~2.2 of PHE-17-a.
\end{proposition}

\begin{proof}
Condition (i): its basic concepts, uniform acceleration, distance and time as continuous magnitudes, have their sense through the idealization of the practices of measuring lengths along a plane and comparing durations. Condition (ii): its evidence was given in life-world experiences. In the \emph{Two New Sciences} the law is supported by balls rolling down a grooved inclined plane, with times measured by weighing the water that flowed from a vessel during the descent \citep[Third Day]{Galileo1638}. Galileo's working notes show that measurement played a substantial role in reaching the law \citep{Drake1978}, and a reconstruction of the experiment has shown that the procedure he describes yields the regularity he claims \citep{Settle1961}.
\end{proof}

\begin{proposition}\label{p:substitution}
The substitution is a step beyond the mathematization, not a part of it.
\end{proposition}

\begin{proof}
The substitution is the claim that the qualities that the method does not mathematize directly, colors, tastes, sounds, exist only in the perceiving body, while figure, number and motion belong to bodies themselves \citep{Galileo1623}. Nothing in the derivation or confirmation of the law of fall requires this claim: Proposition~\ref{p:founded} uses only measured lengths and weights of water, and is indifferent to whether colors belong to bodies. The claim is an interpretation of what the mathematical description is a description of, added to the description.
\end{proof}

By Definition~2.4 of PHE-17-a, forgetting is the taking of idealized objects for the whole of reality. Proposition~\ref{p:substitution} locates the beginning of that forgetting in a philosophical thesis that accompanied Galileo's physics, not in the physics. This is what makes the account compatible with realism. SCI-02 holds that mature, predictively successful theories are approximately true; Galileo's kinematics, as a limiting case of later mechanics, is approximately true of falling bodies, and nothing in Propositions~\ref{p:founded} and~\ref{p:substitution} denies it. What the phenomenological analysis rejects is the further thesis that the world given in experience is an appearance of the world the kinematics describes. Realism about the results requires the first and not the second.

\section{The pattern recurs and is corrected}

Theorem~3.3 of PHE-17-a holds that crises recur wherever a successful method is technized and transmitted. Proposition~\ref{p:substitution} shows the general form of the recurrence: success, followed by an interpretation on which the method's idealized objects are what the science is about and the life-world phenomena that occasioned it are mere data for those objects. Because the interpretation is separable from the results, its correction does not require abandoning them.

The correction of the replication crisis has this structure. SCI-19-b argues that the crisis was itself an episode of self-correction, a change in the rate at which false findings are corrected, produced by reforms that act on the evidential weight of findings and replications. PHE-17-a adds why the forgetting was possible: techniques of significance testing had come to stand in for the question whether an effect could be seen again in the life-world. The two accounts agree. The reforms reconnected the technique to the experiences that confirm it, which is what reactivation in the sense of Definition~2.3 of PHE-17-a requires at the level of practice, and they did so without discarding statistical inference. That correction occurred is consistent with Theorem~3.3 of PHE-17-a, which states a tendency that recurs, not a decline without remedy.

\begin{objection}
Galileo's experiments were idealized and his reported precision is doubtful. If the evidence was not what he claimed, Proposition~\ref{p:founded} fails.
\end{objection}

\begin{reply}
Condition (ii) of Definition~2.2 of PHE-17-a requires that the evidence for the results be finally provided by life-world experiences, not that any particular report be accurate. Whatever the precision of Galileo's own measurements, the law has since been confirmed by measurements of the same kind, and these are life-world experiences.
\end{reply}

\section{Conclusion}

Galileo discovered by mathematizing and concealed by interpreting. Separating the two shows that the priority of the life-world is compatible with the truth of the science founded on it, and that the recurrent crisis is a recurrent misinterpretation, which science can correct without loss of its results.

\begin{thebibliography}{9}\small
\bibitem[Drake(1978)]{Drake1978} Drake, S. (1978). \emph{Galileo at Work: His Scientific Biography}. University of Chicago Press.
\bibitem[Galilei(1623)]{Galileo1623} Galilei, G. (1623). \emph{Il Saggiatore}. Giacomo Mascardi.
\bibitem[Galilei(1638)]{Galileo1638} Galilei, G. (1638). \emph{Discorsi e dimostrazioni matematiche intorno a due nuove scienze}. Elsevier.
\bibitem[Husserl(1954)]{Husserl1954} Husserl, E. (1954). \emph{Die Krisis der europ\"aischen Wissenschaften und die transzendentale Ph\"anomenologie}. Husserliana VI. Ed. W. Biemel. Martinus Nijhoff.
\bibitem[Settle(1961)]{Settle1961} Settle, T.~B. (1961). An experiment in the history of science. \emph{Science}, 133, 19--23.
\end{thebibliography}

\end{document}
